% TOP_PROBLEM: {"id": 3700035, "title": "Finiteness of Newtonian equilibrium points", "category_id": 8, "category": {"id": 8, "name": "analysis", "display_name": "Analysis", "description": "Limits, continuity, calculus, and function theory.", "slug": "analysis", "order_index": 8, "created_at": "2026-07-31T15:26:25.671Z"}, "status": "open"}
% FIRST_POSTED: null
% ENTRY_KIND: statement_only
% Imported from: batch12.tex; UnsolvedMath ID 3700035
% Compile twice: lualatex 3700035.tex
\documentclass[11pt]{article}
\usepackage{fontspec}
\setmainfont{Latin Modern Roman}
\usepackage{amsmath,amssymb,mathtools,unicode-math}
\setmathfont{Latin Modern Math}
\usepackage{xurl,hyperref}
\hypersetup{hidelinks,pdftitle={UnsolvedMath Problem 3700035}}
\newcommand{\sref}[2]{\hyperlink{source:#1:#2}{[S#2]}}
\newcommand{\cataloguescope}[1]{\paragraph{Mathematical question.}#1}
\begin{document}
% UnsolvedMath ID 3700035; source catalogue code AMR-036-0035
\setcounter{section}{3700034}
\section{Finiteness of Newtonian equilibrium points}\label{3700035}
{\small\sffamily UnsolvedMath ID 3700035 · Analysis}
\subsection{Definitions and mathematical statement}

\paragraph{Shared notation.}$\mathbb N=\{1,2,\ldots\}$, and $\mathbb Z,\mathbb Q,\mathbb R,\mathbb C$ denote the integers, rationals, reals and complex numbers. Any other conventions are specified locally.


Let $n\in\mathbb N$, let $x_1,\ldots,x_n\in\mathbb R^3$ be distinct points, and let $a_1,\ldots,a_n>0$. Repeated locations can be combined by adding their positive charges. On $\Omega=\mathbb R^3\setminus\{x_1,\ldots,x_n\}$ define
\[
 u(x)=\sum_{k=1}^{n}\frac{a_k}{\lVert x-x_k\rVert},\qquad
 \operatorname{Crit}(u)=\{x\in\Omega:\nabla u(x)=0\}.
\]
In coordinates this is the solution set of
\[
 \sum_{k=1}^{n}a_k\frac{x-x_k}{\lVert x-x_k\rVert^3}=0.
\]
The locations of the charges are excluded from the critical set.
\paragraph{Statement recorded in the supplied source.}
Is $\operatorname{Crit}(u)$ finite for every such finite positive-charge configuration? No generic-position or nondegeneracy assumption is imposed. In particular, a bound for isolated equilibria alone does not establish this assertion: non-isolated equilibria must also be excluded. This finiteness question is weaker than Maxwell's proposed $(n-1)^2$ bound. \sref{3700035}{2}
\subsection{Short English statement}
Must finitely many positive point charges in three-dimensional space have only finitely many equilibrium points, even for nongeneric configurations?
\subsection{Sources}
\begin{description}
\item[\hypertarget{source:3700035:1}{S1}] ULAM.AI and the UnsolvedMath Contributors, UnsolvedMath: A Curated Collection of Open Mathematics Problems, record 3700035 (AMR-036-0035). CC BY 4.0. \url{https://huggingface.co/datasets/ulamai/UnsolvedMath}.
\item[\hypertarget{source:3700035:2}{S2}] Herbert Edelsbrunner, Christopher Fillmore and Gonçalo Oliveira, Counting equilibria of the electrostatic potential (2025), §1, equation (1), Conjecture 1 and Result 1, pp.1–3. \url{https://arxiv.org/abs/2501.05315}.
\end{description}
\label{end:3700035}
\end{document}
